\chapter{Závěr}
\label{chap:conclusion}

Hlavním cílem této diplomové práce bylo navrhnout, zdokumentovat a naimplementovat interaktivní výukovou aplikaci, která studentům usnadní pochopení datových struktur a třídicích algoritmů.
Tento cíl byl úspěšně splněn vytvořením komplexního webového nástroje \textit{Treeline}, který propojuje teoretické studijní materiály s moderními přístupy v softwarovém inženýrství.

Z inženýrského hlediska je největším přínosem práce návrh modulární a škálovatelné architektury.
Aplikace striktně odděluje aplikační logiku algoritmů od prezentační vrstvy.
Toho bylo dosaženo implementací událostmi řízeného modelu pro asynchronní operace nad datovými strukturami a mechanismem generování předpočítaných stavů pro plynulou animaci třídicích algoritmů.
Úspěšně se podařilo překonat i technologické limity použitých knihoven třetích stran, například vývojem vlastní nezávislé anotační vrstvy nad plátnem \texttt{vis-network} a adaptivním výpočtem překryvů v rekurzivních algoritmech.
Z uživatelského hlediska aplikace nabízí plně responzivní prostředí, dynamické přepínání programovacích jazyků pseudokódu, podporu tmavého režimu a plynulou mezinárodní lokalizaci jak v interaktivních prvcích stránky a jednotlivých vizualizacích, tak v teoretických materiálech, jež vizualizace doplňují.

Zásadní výhodou výsledného řešení je jeho infrastruktura.
Díky využití generování statických stránek (SSG), automatizované CI/CD linky přes GitHub Actions a hostingu prostřednictvím GitHub Pages je aplikace zcela nezávislá na komplexních backendových serverech.
Tento bezúdržbový model provozu v kombinaci s nasazením na vlastní doméně zajišťuje dlouhodobou udržitelnost a dostupnost projektu do budoucnosti s minimálními provozními náklady -- přístup vhodný pro aplikaci zaměřenou na vzdělávání a jednoduchou integraci do výuky.
Aplikace má tak velký potenciál stát se nejen cenným podpůrným nástrojem pro oficiální univerzitní výuku předmětů zaměřených na algoritmy, ale může stabilně sloužit i široké veřejnosti a studentům mimo akademickou půdu.

Navržená architektura představuje pevný základ, který je připraven pro další rozšiřování.
V rámci budoucího vývoje se nabízí implementace dalších teoretických konceptů, pro které je stávající systém již připraven.
Patří mezi ně zejména doména grafových algoritmů (např. procházení grafu do šířky a do hloubky, hledání nejkratší cesty, konstrukce minimální kostry či toky v sítích), komplexnější datové struktury jako hašovací tabulky, nebo různé metody průchodu stromem (inorder, preorder, postorder).

Dalším významným směrem pro budoucí rozvoj je rozšíření aplikace o evaluační modul.
Začlenění interaktivních testů, kvízů a praktických cvičení by studentům umožnilo ověřit si, zda danou problematiku po zhlédnutí vizualizace skutečně plně pochopili, a poskytlo by jim cennou zpětnou vazbu.

Výsledná aplikace dokazuje, že i abstraktní a obtížně uchopitelné koncepty teoretické informatiky lze prezentovat vizuálně atraktivní, moderní a přístupnou formou.
